\section{Algorithms and results, one family at a time}\label{sec:algos}

Each family lives on its own branch (Table~\ref{tab:branches}). Unless stated otherwise, a line extractor is paired
with the benchmark's default depth stage, so differences come from the line alone. Full comparison tables for all
sets are in Section~\ref{sec:results}.

\begin{table}[H]
\centering\small
\caption{Branches. All branch from \texttt{research/base} (simulator + benchmark) and were merged into
\texttt{research/integration} for re-scoring.}\label{tab:branches}
\begin{tabular}{@{}lll@{}}
\toprule
Branch & Stage it changes & Uses learning? \\
\midrule
\texttt{research/peak-detectors} & finding the line (per image column) & no \\
\texttt{research/dp-tracking} & finding the line (global path) & no \\
\texttt{research/cnn-segmentation} & finding the line & yes, U-Net (85\,k weights) \\
\texttt{research/pixel-classifier} & choosing which candidates are laser & yes, gradient boosting \\
\texttt{research/robust-depth} & line $\rightarrow$ groove depth & no \\
\texttt{research/calibration} & pixels $\rightarrow$ millimetres & no \\
\texttt{research/quality-gate} & ``is this photo usable?'' & yes, logistic regression \\
\texttt{research/onoff-fusion} & two photos: laser on, laser off & no \\
\texttt{research/end-to-end} & image $\rightarrow$ depth directly & yes, CNN (155\,k weights) \\
\texttt{research/best-stack} & combination of the above & (gate only) \\
\bottomrule
\end{tabular}
\end{table}

%---------------------------------------------------------------
\subsection{Sub-pixel peak detectors}\label{sec:peak}

\paragraph{Idea.} The laser line runs left to right, so in every image column there is one brightness peak. Find it
and refine its position below one pixel. This is the classic approach in laser scanners and is far cheaper than the
baseline's 2-D ridge detector. We compared centre of gravity, a 3-point Gaussian fit, the Blais--Rioux filters, and
a smoothed peak with a parabola fit; all after removing the background with a morphological top-hat.

\begin{lstlisting}[caption={Smoothed peak + parabola (\texttt{algos/pk\_fir2d.py}). The parabola through the three samples around the maximum gives the sub-pixel row.}]
def parabola3(W, off):                      # W: (N, 3) response samples around each column peak
    h = len(off) // 2
    a, b, c = W[:, h - 1], W[:, h], W[:, h + 1]
    den = a - 2 * b + c
    d = np.where(den < -1e-6, 0.5 * (a - c) / np.where(den < -1e-6, den, -1.0), 0.0)
    return np.clip(d, -1.0, 1.0)

resp = smooth(img, 1.5, 1.5)                # img = top-hat of the frame
thr = max(12.0, 0.25 * np.median(resp.max(axis=0)))
iv, iu = column_peaks(resp, thr, max_per_col=3)
v = iv + parabola3(*windows(resp, iv, iu, 1))
\end{lstlisting}

\paragraph{Results.} In clean conditions every per-column detector beats the baseline: depth MAE on \texttt{ideal}
falls from 0.147 to about 0.06\,mm, and line error from 0.16 to 0.11--0.12\,px. They are also the cheapest methods
(3--5\,ms of their own CPU time against 11\,ms for Steger). Under stress the \emph{typical} error is excellent
(median 0.09\,mm against 0.54\,mm), but the mean stays near 1\,mm, because a per-column detector has no idea of a
\emph{line}: one glint that is brighter than the laser in a few columns becomes a false groove. The 3-point Gaussian
fit collapses when the line saturates (verdict accuracy 0.83 on \texttt{stress}). On real photos only the variant
with an extra pass along the rows (\algo{pk\_fir2d\_rows}, which recovers steep groove walls) finds all six truck
grooves.

\takeaway{Per-column peaks are the most accurate and cheapest way to \emph{locate} a line you have already
identified. They are not enough on their own to decide \emph{which} bright pixels are the line.}

%---------------------------------------------------------------
\subsection{Dynamic-programming line tracking}\label{sec:dp}

\paragraph{Idea.} Use the fact that the laser line is one connected curve crossing the image. Give every pixel a
``reward'' (how line-like it is) and find the single left-to-right path with the best total reward, paying a small
price for gentle slopes and a fixed price for a jump (a groove edge). This is the Viterbi algorithm; it is exact,
needs no training, and one glint cannot win because the path would have to pay to jump to it and back.

\begin{lstlisting}[caption={Core of the Viterbi pass (\texttt{algos/dp\_dense.py}): cost of the best path ending at each pixel, column by column.}]
# pen(d) = A*|d|       for |d| <= s   (gentle slope, e.g. a lit wall: cheap)
#        = J + B*|d|   otherwise      (a groove edge: fixed price + a little per px)
for c in range(1, W):
    full = C[c - 1]                                   # C[c, v] = cost of the best path ending at (c, v)
    prev = full[s:s + H]
    near = prev.copy()
    for k in range(1, s + 1):                         # smooth move of k rows up or down: A per px
        np.minimum(near, full[s - k:s - k + H] + A * k, out=near)
        np.minimum(near, full[s + k:s + k + H] + A * k, out=near)
    # jump from any row v': prev[v'] + J + B*|v - v'|, as two running minima (from above, from below)
    far = J + np.minimum(acc(prev - Br) + Br, acc((prev + Br)[::-1])[::-1] - Br)
    C[c, s:s + H] = np.minimum(near, far) - R[c]      # R = reward in 0..1
\end{lstlisting}

\paragraph{Results.} \algo{dp\_dense} is the best single-frame, training-free line extractor in the study
(Figure~\ref{fig:dp}). On \texttt{stress} the line F1 rises from 0.82 to 0.94, missed grooves fall from 9.9\,\% to
1.8\,\%, depth MAE from 1.52 to 0.54\,mm and P95 from 6.7 to 0.59\,mm, with no unsafe PASS. On \texttt{handheld}
it finds 96\,\% of grooves where the baseline finds 53\,\%. The work on this branch also explained two baseline
faults: Steger's $\pm0.5$\,px acceptance test leaves holes that split one groove into two, and on a saturated line
Steger marks the two \emph{edges} of the bright band instead of its centre. DP takes the mid-point of the band at
half maximum instead. Cost: about 1.4$\times$ the baseline's run time.

\begin{figure}[H]
\centering
\begin{subfigure}[b]{0.54\linewidth}\includegraphics[width=\linewidth]{\figs/dp-tracking/factors_baseline_vs_dp.png}\end{subfigure}\hfill
\begin{subfigure}[b]{0.40\linewidth}\includegraphics[width=\linewidth]{\figs/dp-tracking/real_saturated.png}\end{subfigure}
\caption{Left: baseline (left column) against DP tracking (right column) on the same simulated scene with one
difficulty switched on at a time; yellow bands are reported grooves. Right: two real photos with a thick saturated
line; the baseline (left) follows an edge of the band, DP (right) follows its centre.}
\label{fig:dp}
\end{figure}

\paragraph{What still fails.} Frames with many glints still produce about 0.1 false grooves per frame under stress,
and paint or mud that hides the laser cannot be tracked through. On real photos DP agrees with the reference to
1.1\,px (five of six within 3\,px).

\takeaway{A global ``one line'' constraint removes most gross errors and costs no training data. This is the line
extractor we recommend.}

%---------------------------------------------------------------
\subsection{Learned line detection: a small U-Net}\label{sec:cnn}

\paragraph{Idea.} Train a neural network to paint a heat-map of where the laser line is, then read the peak in each
column. The network is a three-level U-Net with 84{,}705 weights (0.36\,MB), trained only on the 1{,}200 simulated
\texttt{train} frames with a soft cross-entropy plus Dice loss.

\paragraph{Results.} Trained without augmentation, the network gives the best single-frame numbers on simulated
stress (MAE 0.236\,mm, P95 0.47\,mm, verdict 0.985) but fails on real photos: it disagrees with the reference by
7\,px. The cause is that every simulated line sits at the same height in the frame and has the same thickness; the
network memorised that. With augmentation (random shifts, rotation, thickness and gain) real agreement improves
to 1.1\,px and \texttt{handheld} line F1 reaches 0.90, but stress MAE rises to 0.67\,mm. It still needs a hand-written
``keep one line'' rule after the heat-map, and at 60--90\,ms per frame on a laptop core it would take seconds on the
STM32H7.

\begin{figure}[H]
\centering
\includegraphics[width=0.9\linewidth]{\figs/cnn-segmentation/real_refs.png}
\caption{The augmented U-Net on the twelve real reference photos, which it never saw in training
(orange: reference, magenta: network). It follows the line on the truck tyre (bottom row) but is confused by
background clutter on the rover tyre (top row).}
\label{fig:cnn}
\end{figure}

\takeaway{A network trained only on simulation reaches roughly the same real-photo agreement as DP tracking at
more than twice the cost. It becomes worth revisiting once there are a few hundred \emph{real} holder images to train on.}

%---------------------------------------------------------------
\subsection{Candidate classifier}\label{sec:cand}

\paragraph{Idea.} Keep the classical detector, but let a small learned model decide which of its candidate points
are laser and which are glints. Each candidate gets 30 hand-made features (brightness, width, how well it fits its
neighbours, distance from the most common row) and a gradient-boosted classifier (scikit-learn
\texttt{HistGradientBoostingClassifier}, 487\,kB) keeps those with probability $\geq 0.95$.

\paragraph{Results.} On simulated stress this is very strong: MAE 0.273\,mm, false grooves 0.02 per frame, verdict
0.970. But the feature-importance plot (Figure~\ref{fig:cand}) shows it relies almost entirely on one feature, the
distance from the most common row, which works because the simulated line always sits at row $115\pm1$. On real
photos it found two of the six reference grooves. A version without that feature is honest but weaker.

\begin{figure}[H]
\centering
\includegraphics[width=0.85\linewidth]{\figs/pixel-classifier/feature_importance.png}
\caption{What the candidate classifier relies on. Left: the full model leans on ``distance from the most common row''
(\texttt{abs\_dv\_mode}), a shortcut that only exists in the simulator. Right: without it, the model uses local fit
quality instead.}
\label{fig:cand}
\end{figure}

\takeaway{A good score on simulation can come from a shortcut. Looking at what the model uses, and testing on real
photos, exposed it. This result is the reason the \texttt{handheld} and \texttt{shifted} sets were added.}

%---------------------------------------------------------------
\subsection{Robust depth stage}\label{sec:depth}

\paragraph{Idea.} Suppose the line is perfect. How wrong is the depth? With the default stage, about 0.07\,mm on
average and too deep. This branch studied the second half of the pipeline (line $\rightarrow$ depth) on its own and
rebuilt it in five small steps:
\begin{enumerate}
\item Convert every line point to a height and to its true sideways position on the tread, and fit the tread surface
      in millimetres rather than in image columns.
\item Fit the surface with RANSAC (a quadratic). Grooves are one-sided outliers, and the alternatives we tried
      (Huber, Theil--Sen) break when grooves fill much of the view.
\item Detect grooves, including shallow ones, with a running median, and drop narrow shallow blips.
\item Take the groove floor as the deepest cluster found by DBSCAN, and its median as the depth. A single stray
      point can no longer set the depth.
\item Subtract a small tyre-curvature term, and attach an uncertainty and a ``floor not fully visible'' flag.
\end{enumerate}

\begin{lstlisting}[caption={Step 1 (\texttt{algos/robust\_depth.py}): undo lens distortion and express each point's sideways position in mm before fitting the tread surface.}]
def lateral_px(u, v, z, width=320, height=240):
    f = F_PX * width / 320.0
    xd, yd = (u - (width - 1) / 2) / f, (v - (height - 1) / 2) / f
    x, y = xd.copy(), yd.copy()
    for _ in range(3):                         # undo the radial distortion
        k = 1 + K1 * (x * x + y * y)
        x, y = xd / k, yd / k
    return 160.0 + 2.9 * x * (ZC_MM - z)       # X in mm, expressed as "ideal pixels"

# in stage():   u = lateral_px(px, v, -q);  coef = RANSAC quadratic of q against u;  r = q - polyval(coef, u)
# per groove:   z_s = -polyval(coef, u_floor);  z_f = z_s - d_raw
#               d = (z_s - z_f) - TAN_LASER**2 * (z_f**2 - z_s**2) / (2 * tyre_R)
\end{lstlisting}

\paragraph{Results.} With a perfect line, depth MAE on \texttt{ideal} falls from 0.071 to 0.017\,mm and the bias
from $+0.068$ to $-0.012$\,mm (Figure~\ref{fig:depth}, top). Most of the old bias was \emph{not} tyre curvature, as
we first believed: it came from camera parallax and lens distortion bending the crown in image columns. The
curvature term itself only matters for deep grooves ($+0.09$\,mm at 17\,mm) and is insensitive to the assumed tyre
radius ($\pm0.005$\,mm for $\pm20\,\%$).

With the baseline line, the new stage (\algo{rd\_final}) cuts stress MAE from 1.52 to 0.46\,mm and, on worn tyres,
MAE from 0.053 to 0.019\,mm with \emph{no missed grooves}. The shallow-groove rule matters for safety: on a
special set of 300 unevenly worn tyres the default stage called 9\,\% of illegal tyres PASS because it did not see
the shallowest groove at all; the new stage called none (Figure~\ref{fig:depth}, bottom).

\begin{figure}[H]
\centering
\includegraphics[width=0.92\linewidth]{\figs/robust-depth/bias_vs_depth.png}\\[4pt]
\includegraphics[width=0.74\linewidth]{\figs/robust-depth/legal_boundary.png}
\caption{Top: depth error with a perfect line against true depth: default stage, surface fitted in millimetres, and
with the curvature term added. Bottom: measured against true depth on unevenly worn tyres; red crosses are grooves
the stage did not see. The vertical and horizontal lines mark 1.6\,mm.}
\label{fig:depth}
\end{figure}

\paragraph{Limits.} Rocking the holder about the tread axis adds about 1.2\,\% of the depth per degree, and one
laser line cannot detect it. The correction uses holder constants (camera height, lens distortion), and the
simulated crown is an exact parabola, so the gain on real treads will be smaller than in simulation. On real
photos the DBSCAN floor disagrees with two of the six references by about 2.4\,px; without a gauge reading we
cannot say which is right, but on the one pair of repeated photos it is far more consistent (0.17\,px apart instead
of 2.45\,px).

\takeaway{The depth stage was the largest source of error after gross line mistakes, and fixing it needs no
learning: fit in millimetres, take the floor from a cluster, look for shallow grooves explicitly.}

%---------------------------------------------------------------
\subsection{Calibration}\label{sec:calib}

\paragraph{Idea.} Calibration turns ``the line is at pixel $(u,v)$'' into ``the surface is at height $z$''. The
October procedure photographs a flat plate at known heights and fits a polynomial. We compared polynomials, splines,
a Gaussian process, and a \emph{physical} model with 5 or 7 parameters (pinhole camera, one distortion term, and
the laser plane).

\begin{figure}[H]
\centering
\includegraphics[width=0.9\linewidth]{\figs/calibration/a_plates.png}\\[4pt]
\includegraphics[width=0.9\linewidth]{\figs/calibration/d_drift.png}
\caption{Top: calibration error against the number of plate heights (left) and their range (right); the physical
models (red, orange) stay accurate with three plates. Bottom: depth error after the laser is knocked by a fraction
of a degree (left), and after correcting with one photo of a flat plate at zero height (centre, right).}
\label{fig:calib}
\end{figure}

\paragraph{Results.}
\begin{itemize}
\item With 15 plates every sensible model is accurate to better than 0.02\,mm. With only three plates the physical
      model still gives 0.002\,mm; the polynomial gives 0.12\,mm.
\item The October recipe (plates from 0 to $-20$\,mm) never covers the tread itself, which sits at $+3.2$\,mm, so the
      polynomial has to extrapolate. Recommended heights: $+4$, $+2$, $0$, $-3$, $-6$, $-10$, $-14$, $-18$\,mm.
\item If the laser is knocked by $0.5^\circ$, every depth is scaled by about $+4\,\%$ (0.67\,mm at 16\,mm). One photo
      of a flat plate at zero height detects this, and re-aiming the laser plane in the physical model corrects it
      to within 0.003\,mm. A change of focal length is \emph{not} visible to that check.
\item A systematic error in the shim heights passes straight through every model. Only a step of known depth can
      catch it.
\end{itemize}

\caution{The simulator's camera \emph{is} a pinhole with one distortion term, so the physical model fits it perfectly
by construction. Expect a smaller advantage on the real device; the ranking for few plates and for extrapolation
should still hold.}

\takeaway{Use the physical model, calibrate with eight heights that bracket the tread, and add a daily
``zero check'' on a flat plate plus a known-depth step block.}

%---------------------------------------------------------------
\subsection{Quality gate}\label{sec:gate}

\paragraph{Idea.} The baseline reports a depth on about 70\,\% of photos that a person would refuse to measure. A
\emph{gate} decides, before a number is shown, whether the photo is good enough. We built it in two stages: a
\emph{frame gate} (is there a laser line on a tyre at all?) trained on the 86 labelled real photos, and a \emph{risk
gate} (is this particular reading likely to be wrong?) trained on simulated frames where the error is known.
Both are regularised logistic regressions on a few dozen hand-made features.

\begin{lstlisting}[caption={Some of the risk-gate features (\texttt{algos/gate\_features.py}). Each is a number a technician could also judge by eye.}]
r = v - np.polyval(coef, u)                       # > 0 = below the fitted tread surface
z.update({
    "min_floor_pts": float(min(floor_n, default=0)),          # fewest points on any groove floor
    "depth_spread_px": float(np.ptp([x["depth"] for x in g])) if g else 0.0,   # grooves of one tyre wear alike
    "above_frac": float((r < -depth.PX["min_depth"]).mean()),  # points ABOVE the tread: glints, never rubber
    "max_gap_frac": float(np.max(np.diff(np.r_[0, uc, w - 1])) / w),           # biggest hole in the line
    "multi_frac": float((vr > 3).mean()),                      # columns where the line has two branches
})
\end{lstlisting}

\paragraph{Results.} With the gate, the baseline answers 88\,\% of \texttt{ideal} frames, 96\,\% of \texttt{legal}
frames and only 16\,\% of \texttt{stress} frames, but the answers it does give under stress are good: MAE falls from
1.52 to 0.14\,mm and gross errors from 36\,\% to zero (Figure~\ref{fig:gate}). On real photos it keeps 17 of the 19
usable frames and cuts false measurements from 69\,\% to 4\,\% on frames with no laser, and from 75\,\% to 25\,\% on
frames where the laser is visible but unusable.

\begin{figure}[H]
\centering
\includegraphics[width=0.8\linewidth]{\figs/quality-gate/risk_coverage.png}
\caption{Risk--coverage curves: if the device only answers the share of presses it trusts most (horizontal axis),
how bad are the answers (vertical axis)? Dots mark the threshold shipped, fixed on \texttt{train}.}
\label{fig:gate}
\end{figure}

\caution{The frame gate was trained on the same 86 real photos, so the real-photo numbers above are
\emph{in-sample}. The honest estimate is leave-one-session-out cross-validation: a false-accept rate of 0.24
(95\,\% interval 0.14--0.36) when every choice is nested inside the cross-validation. More complex models did not
help: a small CNN gate was at chance on a new day (AUC 0.49).}

\takeaway{``Retake the photo'' is the cheapest robustness feature. With only 86 labelled photos, simple logistic
regression is the right-sized model; more labelled sessions will help more than a bigger model.}

%---------------------------------------------------------------
\subsection{Laser on/off subtraction}\label{sec:onoff}

\paragraph{Idea.} Take two photos a few milliseconds apart, one with the laser on and one with it off, and
subtract. Everything that is not laser (sunlight, glints, paint) cancels. The October photos have no laser-off
partner, so this was tested in simulation only, as a question to the hardware team: is it worth letting the
microcontroller switch the laser?

\begin{figure}[H]
\centering
\includegraphics[width=\linewidth]{\figs/onoff-fusion/example_on_off_diff.png}\\[4pt]
\includegraphics[width=0.82\linewidth]{\figs/onoff-fusion/error_vs_motion.png}
\caption{Top: a stress frame with glints; laser on, laser off, their difference, the difference when the hand moved
2\,px between the photos, and the same after registration. Bottom: error against hand motion between the two
photos, for plain and registered subtraction.}
\label{fig:onoff}
\end{figure}

\paragraph{Results.} With a simple centre-of-gravity detector on the difference image (\algo{onoff\_cog}), stress
MAE is 0.147\,mm, P95 0.40\,mm, false grooves 0.01 per frame, and no unsafe PASS on that set: the best
stress numbers of any single method. Motion matters: with 2\,px of movement between the photos, plain
subtraction is \emph{worse} than a single frame. Matching the gain, registering the two frames by phase correlation
and dilating the off frame by one pixel keeps the error flat up to 3\,px of motion. Subtraction does not help with a
saturating background, rocking, or a dim laser. Averaging several presses helps little, because the remaining
error is mostly bias; a median of three is cheap insurance.

\takeaway{If the next board revision can switch the laser with a GPIO pin, on/off subtraction is the single biggest
robustness gain available, and it is twenty lines of code.}

%---------------------------------------------------------------
\subsection{End-to-end learning}\label{sec:e2e}

\paragraph{Idea.} Skip the geometry: train a network to output groove depths directly from the image. A CNN with
154{,}678 weights predicts, for every image column, the probability of a groove and its depth.

\paragraph{Results.} On the sets that look like its training data it does well (\texttt{ideal} MAE 0.100\,mm,
\texttt{stress} 0.483\,mm, 2\,ms per frame). Then it fails in every way that matters for a measuring instrument
(Figure~\ref{fig:e2e}):
\begin{itemize}
\item On another holder unit (\texttt{shifted}) the MAE is 5.2\,mm with 8.7\,\% unsafe PASS. A calibration cannot
      fix it, because the network has no calibration input.
\item With the laser switched \emph{off} it still ``measures'' 52--91\,\% of the grooves. It learned that grooves look
      dark, and that darker means deeper in the simulator.
\item On 67 real photos that cannot be measured it reported a depth on 43, and never said FAIL.
\item At the legal boundary it is the least safe method tested (2.7\,\% unsafe PASS, against 0 for the baseline).
\end{itemize}

\begin{figure}[H]
\centering
\includegraphics[width=0.56\linewidth]{\figs/end-to-end/real_and_laser_off.png}\hfill
\includegraphics[width=0.42\linewidth]{\figs/end-to-end/ood_error.png}
\caption{Left: the end-to-end network reports confident depths on photos of people and a ceiling light, and still
finds grooves with the laser off. Right: its error under changes it was not trained for, against the geometric
baseline.}
\label{fig:e2e}
\end{figure}

\takeaway{For a safety measurement, keep the geometry as the measurement. Triangulation is correct for a reason we
can explain to an auditor; a network's answer is only correct while the pictures look like its training set. Use
learning as a helper (gate, line finder, cross-check), not as the ruler.}
